\documentclass[11pt]{amsart}
\usepackage{amsmath,amssymb,amsthm,booktabs}
\usepackage[margin=1.15in]{geometry}
\usepackage[hidelinks]{hyperref}
\newtheorem{theorem}{Theorem}[section]
\newtheorem{lemma}[theorem]{Lemma}
\newtheorem{proposition}[theorem]{Proposition}
\newtheorem{corollary}[theorem]{Corollary}
\theoremstyle{remark}
\newtheorem{remark}[theorem]{Remark}
\newcommand{\Z}{\mathbb{Z}}
\newcommand{\cU}{\mathcal{U}}
\newcommand{\BP}{\mathrm{BP}}
\newcommand{\doi}[1]{\href{https://doi.org/#1}{doi:\nolinkurl{#1}}}
\newcommand{\arxiv}[1]{\href{https://arxiv.org/abs/#1}{arXiv:\nolinkurl{#1}}}
\setlength{\emergencystretch}{2.5em}

\title[Stationary measures of planar kinetically constrained models]{Stationary measures of planar kinetically
constrained models with bounded conditional odds}
\author{Yitzchak Shmalo}
\address{Einstein Institute of Mathematics, The Hebrew University of Jerusalem, Jerusalem, Israel}
\email{yspennstate@gmail.com}
\date{November 2026}
\subjclass[2020]{60K35, 82C22}
\keywords{kinetically constrained models, Fredrickson--Andersen model, stationary measures, entropy production,
bootstrap percolation}

\raggedbottom
\begin{document}
\begin{abstract}
Martinelli, Shapira and Toninelli conjectured that every stationary measure of a kinetically constrained spin model
on $\Z^d$ is a mixture of the equilibrium product measure conditioned on the fibers of bootstrap percolation. In
dimension one, and for translation invariant stationary measures in every dimension, stationarity is known to imply detailed
balance. We prove that on $\Z^2$ every stationary measure whose single-site conditional probabilities are bounded away
from $0$ and $1$ satisfies detailed balance, without translation invariance and although the transition rates vanish
on a large set. If large boxes are also internally spanned at every positive density, such a measure is the
equilibrium product measure; this applies to the one- and two-neighbour Fredrickson--Andersen models at every density.
The proof adapts the entropy argument of Holley and Stroock to degenerate rates, the hypothesis on the conditional
probabilities replacing a lower bound on the rates. We also record the classification of all stationary measures of
the one-neighbour model on strips $\Z\times H$ with $H$ finite.
\end{abstract}
\maketitle

\section{Introduction}\label{sec:intro}

We recall the models and the conjecture of Martinelli, Shapira and Toninelli, and state the results.

\subsection{The models}
Let $d\ge1$, let $G$ be $\Z^d$ or a strip $\Z\times H$ with $H$ a finite connected graph, and let
$\Omega=\{0,1\}^G$. A site $x$ with $\eta_x=0$ is \emph{empty} and a site with $\eta_x=1$ is \emph{filled}. Fix
$q\in(0,1)$, put $p=1-q$, and let $\pi=\pi_q$ be the product measure on $\Omega$ with $\pi(\eta_x=0)=q$ for every
$x$; for $a\in(0,1)$, $\pi_a$ denotes the product measure with $\pi_a(\eta_x=0)=a$. To every site $x$ is attached a
\emph{constraint} $c_x\colon\Omega\to\{0,1\}$ that depends only on the spins at finitely many sites other than $x$.
We assume a uniform range: there is $r\ge1$ such that $c_x$ depends only on spins at sites $y\ne x$ whose first
coordinates differ from that of $x$ by at most $r$ (on $\Z^2$: whose $\ell^\infty$ distance from $x$ is at most $r$).
The kinetically constrained model is the Markov process on $\Omega$ with generator
\begin{equation}\label{eq:generator}
Lf(\eta)=\sum_{x\in G}c_x(\eta)\bigl(\pi_x(f)-f\bigr)(\eta)=\sum_{x\in G}c_x(\eta)\,r_x(\eta)\bigl(f(\eta^x)-f(\eta)\bigr),
\end{equation}
where $\pi_x(f)$ averages $f$ over the spin at $x$ with respect to $\pi$, $\eta^x$ is $\eta$ with the spin at $x$
flipped, and $r_x(\eta)=q$ if $\eta_x=1$, $r_x(\eta)=p$ if $\eta_x=0$. Thus a filled site empties at rate $q$ and an
empty site fills at rate $p$, provided its constraint is satisfied. The rates are bounded and of finite range, so $L$
generates a Feller semigroup $(P_t)$, and $\pi$ is reversible \cite{CMRT08,HT25}. A probability measure $\mu$ on
$\Omega$ is \emph{stationary} if $\mu P_t=\mu$ for all $t\ge0$; then $\mu(Lf)=0$ for every local function $f$, that is
every function of finitely many spins. No translation invariance, reversibility or absolute continuity is assumed.

The main examples are given by an \emph{update family} $\cU$, a finite collection of finite nonempty subsets of
$\Z^d\setminus\{0\}$, through
\begin{equation}\label{eq:update}
c_x(\eta)=\mathbb I\{\text{there is }U\in\cU\text{ with }\eta_{x+u}=0\text{ for all }u\in U\},
\end{equation}
where $\mathbb I$ denotes an indicator.
Such a constraint is \emph{monotone}: emptying a site never turns a satisfied constraint into an unsatisfied one. The
Fredrickson--Andersen $j$-neighbour model (FA-$j$f) \cite{FA84} has $c_x=1$ when at least $j$ nearest neighbours of $x$
are empty. On a strip the same definitions apply, with constraints that may depend on the site.

\emph{Bootstrap percolation} empties, one after another, filled sites whose constraint is satisfied and never fills a
site. For monotone constraints the limit $\BP(\eta)$ does not depend on the order. For a finite $V\subset G$ let
$\mathcal G_V$ be the event that bootstrap percolation run inside $V$ from $\eta_V$, with every site outside $V$
filled, empties all of $V$; we then say that $V$ is \emph{internally spanned}. The event $\mathcal G_V$ depends only on
$\eta_V$ and is increasing in the set of empty sites.

\subsection{The conjecture and what is known}
Let $E=\{\beta\in\Omega:\BP(\beta)=\beta\}$. Martinelli, Shapira and Toninelli \cite[Conjecture~3]{MST25} conjectured
that every stationary measure $\mu$ has the form $\mu=\int_E\pi(\,\cdot\mid\BP=\beta)\,\mu^*(d\beta)$ for a probability
$\mu^*$ on $E$, the mixing measure; see also \cite[Conjecture~1.6]{Mar26}. For the conditioning at fibers of $\pi$-probability zero a
specific choice of kernel is part of the statement; the results below involve only the fibers $\{\BP=0\}$ and, on
strips, the single configuration $\mathbf 1$ with every site filled.

The route in \cite{MST25} has two steps: show that a stationary measure satisfies \emph{detailed balance}, then identify
reversible measures by internal spanning. A probability $\mu$ satisfies detailed balance if for every finite $W\subset
G$ and every $x\in W$ such that $c_x$ depends only on spins in $W$,
\begin{equation}\label{eq:db}
c_x(\zeta)\,r_x(\zeta)\,\mu_W(\zeta)=c_x(\zeta^x)\,r_x(\zeta^x)\,\mu_W(\zeta^x)\qquad\text{for all }\zeta\in\{0,1\}^W,
\end{equation}
where $\mu_W$ is the marginal of $\mu$ on $W$. For translation invariant $\mu$ the first step follows from Holley's free
energy argument \cite{Hol71}, \cite[Theorem~2.4]{MST25}; in dimension one it holds without translation invariance
\cite[Theorem~2.5]{MST25}. For FA-1f, BABP and $\delta$-West on $\Z$ this gives the full classification
\cite[Corollary~2.9]{MST25}, and the one-dimensional East model is classified in \cite[Theorem~2.12]{MST25}, the
$d$-dimensional East model in \cite[Theorem~2.13]{MST25} by a different, oriented argument. In dimension two the
entropy argument without translation invariance loses a logarithmic factor at the boundary \cite[Remark~2.16]{MST25}.

For spin systems with rates bounded away from zero, Holley and Stroock \cite{HS77} showed that in dimensions one and
two every stationary measure of a stochastic Ising model is a Gibbs state, without translation invariance; Moulin
Ollagnier and Pinchon \cite{MOP77} gave the free energy argument for translation invariant stationary measures, and Jahnel and
K\"oppl \cite{JK25} proved a general version for interacting particle systems in dimensions one and two under rates
uniformly bounded away from zero. In kinetically constrained models the rates vanish on a set of positive measure, so
these results do not apply directly. For FA-1f at high density, convergence to $\pi$ from initial laws with good control
of the distance to the nearest empty site is known \cite{BCMRT13,MV19,HT25}.

\subsection{Results}
A probability $\mu$ on $\Omega$ has \emph{bounded conditional odds} if there is $\varepsilon\in(0,1/2]$ such that for
every finite $V\subset G$, every $x\in V$ and every $\xi\in\{0,1\}^{V\setminus\{x\}}$ with $\mu(\eta_{V\setminus\{x\}}=\xi)>0$,
\begin{equation}\label{eq:odds}
\varepsilon\le\mu\bigl(\eta_x=0\bigm|\eta_{V\setminus\{x\}}=\xi\bigr)\le1-\varepsilon.
\end{equation}
Revealing the spins of a finite set one at a time shows that every finite configuration then has positive probability.
Absolute continuity of the finite marginals with respect to $\pi$ is no restriction, since $\pi$ charges every finite
configuration; \eqref{eq:odds} is a genuine hypothesis.

\begin{theorem}\label{thm:balance}
Let $G=\Z^2$ and let the constraints be of uniform finite range and independent of the spin they constrain. Every
stationary measure with bounded conditional odds satisfies detailed balance \eqref{eq:db}.
\end{theorem}

\begin{theorem}\label{thm:identify}
Let $G=\Z^2$, let the constraints be given by an update family as in \eqref{eq:update}, and let $V_n=[-n,n]^2$.
Assume that $\pi_a(\mathcal G_{V_n})\to1$ as $n\to\infty$ for every $a\in(0,1)$. Then every stationary measure with
bounded conditional odds equals $\pi$. This applies to FA-1f and FA-2f on $\Z^2$, for every $q\in(0,1)$.
\end{theorem}

For FA-1f the configuration $\mathbf 1$ is frozen and $\delta_{\mathbf 1}$ is stationary, so a stationary measure
need not have bounded conditional odds. Theorem~\ref{thm:identify} applies to the part of $\mu$ off $\mathbf 1$:

\begin{corollary}\label{cor:fa1f}
Let $\mu$ be a stationary measure of FA-1f on $\Z^2$ with $\mu(\{\mathbf 1\})<1$. If
$\nu=(\mu-\mu(\{\mathbf 1\})\delta_{\mathbf 1})/(1-\mu(\{\mathbf 1\}))$ has bounded conditional odds, then
$\mu=(1-\lambda)\pi+\lambda\delta_{\mathbf 1}$ with $\lambda=\mu(\{\mathbf 1\})$.
\end{corollary}

On strips the boundary of a box has bounded size and no hypothesis on $\mu$ is needed.

\begin{proposition}\label{prop:strip}
Let $G=\Z\times H$ with $H$ a finite connected graph, and let $\mu$ be a stationary measure of FA-1f on $G$. Then
$\mu=(1-\lambda)\pi+\lambda\delta_{\mathbf 1}$ with $\lambda=\mu(\{\mathbf 1\})$. The same holds for all monotone
constraints of uniform finite range, possibly depending on the site, that allow every FA-1f move and vanish at
$\mathbf 1$.
\end{proposition}

Proposition~\ref{prop:strip} also follows from \cite[Theorem~2.5 and Lemma~2.8]{MST25}: writing
$H=\{h_0,\dots,h_{m-1}\}$ and sending $(k,h_i)$ to $mk+i$ turns the model into a one-dimensional model with
site-dependent constraints of range at most $mr+m-1$, which those results allow. We include the short direct proof
because it is the case $|B_n|$ bounded of the argument below.

Theorems~\ref{thm:balance} and~\ref{thm:identify} do not settle the conjecture on $\Z^2$: a stationary measure need
not have bounded conditional odds, and on fibers other than $\{\BP=0\}$ the identification of reversible measures is a
separate problem. The proof uses the divergence of $\sum_n1/n$ and does not extend to $d\ge3$
(Remark~\ref{rem:dim}).

Section~\ref{sec:entropy} sets up the entropy balance for an arbitrary stationary measure. Section~\ref{sec:proofs}
proves Theorem~\ref{thm:balance} and the balance part of Proposition~\ref{prop:strip}, Section~\ref{sec:ident} the
identification, and Section~\ref{sec:fa2f} the internal spanning property of FA-2f.

\section{Entropy balance}\label{sec:entropy}

Throughout this section $\mu$ is a stationary measure and the constraints satisfy the hypotheses of
Theorem~\ref{thm:balance}, on $G=\Z^2$ or on a strip. For a finite $V\subset G$, $\zeta\in\{0,1\}^V$ and $x\in V$ put
\[
a_V(x,\zeta)=\int\mathbb I\{\eta_V=\zeta\}\,c_x(\eta)\,r_x(\eta)\,\mu(d\eta),\qquad
b_V(x)=\sum_{\zeta}\bigl|a_V(x,\zeta)-a_V(x,\zeta^x)\bigr|.
\]
Call $x$ \emph{interior} to $V$ if $c_x$ depends only on spins in $V$, and for interior $x$ put
\[
A_V(x)=\sum_\zeta\bigl(a_V(x,\zeta)-a_V(x,\zeta^x)\bigr)\log\frac{a_V(x,\zeta)}{a_V(x,\zeta^x)},
\]
with terms in which both values vanish equal to zero. We write $a=a_V(x,\zeta)$ and $a^x=a_V(x,\zeta^x)$ when $V$ and
$x$ are fixed. Since $c_xr_x\le1$, $\sum_\zeta a_V(x,\zeta)\le1$ and $b_V(x)\le2$.

\begin{lemma}\label{lem:zero}
If $\mu_V(\zeta)=0$ then $a_V(x,\zeta)=a_V(x,\zeta^x)=0$ for every $x\in V$.
\end{lemma}

\begin{proof}
Clearly $a_V(x,\zeta)=0$. Apply $\mu(Lf)=0$ to $f=\mathbb I\{\eta_V=\zeta\}$. A flip outside $V$ does not change
$\eta_V$, so
\[
0=\sum_{y\in V}\sum_{\xi}a_V(y,\xi)\bigl(\mathbb I\{\xi^y=\zeta\}-\mathbb I\{\xi=\zeta\}\bigr)
=\sum_{y\in V}a_V(y,\zeta^y)-\sum_{y\in V}a_V(y,\zeta).
\]
The second sum vanishes, and the first is a sum of nonnegative terms, so $a_V(y,\zeta^y)=0$ for every $y\in V$.
\end{proof}

In particular, at an interior $x$ with $a>0$ and $a^x>0$ both $\zeta$ and $\zeta^x$ have positive mass, and every
pair in the sums below with a vanishing marginal mass contributes zero. Put $\rho_V=\mu_V/\pi_V$ on
$\{\mu_V>0\}$ and
\[
\ell_V(x,\zeta)=\log\frac{\rho_V(\zeta)}{\rho_V(\zeta^x)}\qquad\text{when }\mu_V(\zeta)\mu_V(\zeta^x)>0 .
\]

\begin{lemma}\label{lem:abc}
Let $V\subset W$ be finite.
\begin{enumerate}
\item[(a)] If $x$ is interior to $V$, then $0\le A_V(x)\le A_W(x)$.
\item[(b)] $b_V(x)\le b_W(x)$ for every $x\in V$.
\item[(c)] If $x$ is interior to $V$, then $b_V(x)^2\le A_V(x)$.
\end{enumerate}
\end{lemma}

\begin{proof}
The function $\phi(u,v)=(u-v)\log(u/v)$ on $[0,\infty)^2$, with $\phi(0,0)=0$ and $\phi=+\infty$ when exactly one
argument vanishes, is nonnegative, convex and positively homogeneous of degree one, hence subadditive. Since
$a_V(x,\zeta)=\sum_\xi a_W(x,\xi)$ and $a_V(x,\zeta^x)=\sum_\xi a_W(x,\xi^x)$, the sums running over the extensions
$\xi$ of $\zeta$ to $W$, subadditivity gives (a), and the triangle inequality gives (b). For (c), the logarithmic mean
is at most the arithmetic mean, so $\phi(u,v)\ge2(u-v)^2/(u+v)$. By the Cauchy--Schwarz inequality,
\[
b_V(x)^2\le\sum_\zeta\frac{(a-a^x)^2}{a+a^x}\ \sum_\zeta(a+a^x)\le\frac{A_V(x)}2\cdot2 .\qedhere
\]
\end{proof}

Now fix finite sets $I\subset V$ such that every $x\in I$ is interior to $V$, and put $B=V\setminus I$.

\begin{lemma}\label{lem:balance}
With $M=\max(p/q,q/p)$,
\begin{equation}\label{eq:balance}
\sum_{x\in I}A_V(x)=-\sum_{x\in B}\sum_\zeta\bigl(a_V(x,\zeta)-a_V(x,\zeta^x)\bigr)\ell_V(x,\zeta),
\end{equation}
where the inner sum runs over $\zeta$ with $\mu_V(\zeta)\mu_V(\zeta^x)>0$. For every $\lambda\ge1$ and $x\in B$,
\begin{equation}\label{eq:tail}
-\sum_\zeta\bigl(a_V(x,\zeta)-a_V(x,\zeta^x)\bigr)\ell_V(x,\zeta)\le\lambda\,b_V(x)+2M\lambda e^{-\lambda},
\end{equation}
and the left side of \eqref{eq:tail} is at most $2M/e$.
\end{lemma}

\begin{proof}
Let $F$ be a function on $\{0,1\}^V$ equal to $\log\rho_V$ where $\mu_V>0$ and to $0$ elsewhere, and apply
$\mu(Lf)=0$ to $f(\eta)=F(\eta_V)$:
\[
0=\sum_{x\in V}\sum_\zeta a_V(x,\zeta)\bigl(F(\zeta^x)-F(\zeta)\bigr)
=-\frac12\sum_{x\in V}\sum_\zeta\bigl(a_V(x,\zeta)-a_V(x,\zeta^x)\bigr)\bigl(F(\zeta)-F(\zeta^x)\bigr),
\]
where the second form comes from replacing $\zeta$ by $\zeta^x$ and averaging. By Lemma~\ref{lem:zero} only pairs with
$\mu_V(\zeta)\mu_V(\zeta^x)>0$ contribute, and for them $F(\zeta)-F(\zeta^x)=\ell_V(x,\zeta)$. Let $x\in I$. Then
$c_x$ is a function of $\zeta$ that does not depend on $\zeta_x$, so $a=c_x(\zeta)r_x(\zeta)\mu_V(\zeta)$ and
$a^x=c_x(\zeta)r_x(\zeta^x)\mu_V(\zeta^x)$. Since $r_x(\zeta)/r_x(\zeta^x)=\pi_V(\zeta^x)/\pi_V(\zeta)$, whenever
$c_x(\zeta)=1$ we get $a/a^x=\rho_V(\zeta)/\rho_V(\zeta^x)$, that is $\ell_V(x,\zeta)=\log(a/a^x)$, and when
$c_x(\zeta)=0$ both $a$ and $a^x$ vanish. Hence the terms with $x\in I$ add up to $-\frac12\sum_{x\in I}A_V(x)$, which
proves \eqref{eq:balance}.

For \eqref{eq:tail}, the summand is unchanged when $\zeta$ is replaced by $\zeta^x$, because both factors change
sign. Split the $\zeta$ with $\ell=\ell_V(x,\zeta)$ into $|\ell|\le\lambda$, $\ell>\lambda$ and $\ell<-\lambda$. The
first group contributes at most $\lambda b_V(x)$. If $\ell>\lambda$, then $-(a-a^x)\ell\le a^x\ell\le\mu_V(\zeta^x)\ell$,
and
\[
\mu_V(\zeta^x)=\rho_V(\zeta^x)\pi_V(\zeta^x)=\mu_V(\zeta)\,\frac{\pi_V(\zeta^x)}{\pi_V(\zeta)}\,e^{-\ell}
\le M\mu_V(\zeta)e^{-\ell}.
\]
Since $t\mapsto te^{-t}$ decreases on $[1,\infty)$, the second group contributes at most
$M\lambda e^{-\lambda}\sum_\zeta\mu_V(\zeta)=M\lambda e^{-\lambda}$, and by the symmetry just noted the third
group contributes the same. The same argument with threshold $0$, using $\sup_{t\ge0}te^{-t}=1/e$, bounds the whole
sum by $2M/e$.
\end{proof}

Now let $V_n$ be boxes increasing in $n$, $I_n=V_{n-1}$ and $B_n=V_n\setminus V_{n-1}$, chosen so that every site of
$V_{n-1}$ is interior to $V_n$. Put
\[
S_n=\sum_{x\in I_n}A_{V_n}(x).
\]

\begin{lemma}\label{lem:shell}
$S_n$ is nondecreasing, $0\le S_n\le(2M/e)|B_n|$, and
\begin{equation}\label{eq:current}
\sum_{x\in B_n}b_{V_n}(x)\le\sqrt{|B_n|\,(S_{n+1}-S_n)} .
\end{equation}
\end{lemma}

\begin{proof}
The bound on $S_n$ is \eqref{eq:balance} with the last statement of Lemma~\ref{lem:balance}. Since
$I_{n+1}=I_n\cup B_n$, Lemma~\ref{lem:abc}(a) gives
$S_{n+1}\ge\sum_{x\in I_n}A_{V_n}(x)+\sum_{x\in B_n}A_{V_{n+1}}(x)=S_n+\sum_{x\in B_n}A_{V_{n+1}}(x)$. Every $x\in B_n$
is interior to $V_{n+1}$, so by Lemma~\ref{lem:abc}(b) and (c) and the Cauchy--Schwarz inequality,
\[
\sum_{x\in B_n}b_{V_n}(x)\le\sum_{x\in B_n}\sqrt{A_{V_{n+1}}(x)}\le\sqrt{|B_n|\sum_{x\in B_n}A_{V_{n+1}}(x)}
\le\sqrt{|B_n|(S_{n+1}-S_n)} .\qedhere
\]
\end{proof}

\begin{lemma}\label{lem:zeroS}
If $S_n=0$ for every $n$, then $\mu$ satisfies detailed balance \eqref{eq:db}.
\end{lemma}

\begin{proof}
Every summand of $A_{V_n}(x)$ is nonnegative, so $a_{V_n}(x,\zeta)=a_{V_n}(x,\zeta^x)$ for all $x\in V_{n-1}$ and all
$\zeta$. Let $W$ be finite and $x\in W$ with $c_x$ depending only on spins in $W$. Choose $n$ with $W\cup\{x\}\subset
V_{n-1}$ and sum the identity over the extensions of $\xi\in\{0,1\}^W$ to $V_n$; this gives
$a_W(x,\xi)=a_W(x,\xi^x)$, which is \eqref{eq:db} because $a_W(x,\xi)=c_x(\xi)r_x(\xi)\mu_W(\xi)$ and
$c_x(\xi)=c_x(\xi^x)$.
\end{proof}

\section{Proof of Theorem~\ref{thm:balance}}\label{sec:proofs}

\begin{proof}[Proof of Theorem~\ref{thm:balance}]
Let $V_n=[-nr,nr]^2$, so that every site of $V_{n-1}$ is interior to $V_n$ and
\begin{equation}\label{eq:shellsize}
|B_n|=(2nr+1)^2-(2(n-1)r+1)^2=8r^2n-4r^2+4r\le8r^2n .
\end{equation}
By \eqref{eq:odds}, for every finite $V$, $x\in V$ and $\zeta$, $\mu_V(\zeta)/\mu_V(\zeta^x)$ lies between
$\varepsilon/(1-\varepsilon)$ and $(1-\varepsilon)/\varepsilon$; dividing by $\pi_V(\zeta)/\pi_V(\zeta^x)\in\{p/q,q/p\}$,
\[
|\ell_V(x,\zeta)|\le K:=\log\frac{1-\varepsilon}{\varepsilon}+\Bigl|\log\frac pq\Bigr| .
\]
Then \eqref{eq:balance} and \eqref{eq:current} give
\begin{equation}\label{eq:planar}
S_n\le K\sum_{x\in B_n}b_{V_n}(x)\le K\sqrt{8r^2n\,(S_{n+1}-S_n)},
\end{equation}
and also $S_n\le 2K|B_n|\le16Kr^2n$, since $b_{V_n}(x)\le2$. If $K=0$ then $S_n=0$. Otherwise put $C=8K^2r^2$ and
$D=16Kr^2$, so that
\[
S_n^2\le Cn\,(S_{n+1}-S_n),\qquad S_n\le Dn\qquad(n\ge1).
\]
Suppose $S_m=s>0$ for some $m$. For $n\ge m$, since $S$ is nondecreasing, $S_{n+1}\ge S_n(1+S_n/(Cn))\ge
S_n(1+s/(Cn))$, and $\prod_{n\ge m}(1+s/(Cn))=\infty$; hence $S_n\to\infty$. Choose $m'\ge m$ with $S_n\ge3C$ for
$n\ge m'$. Then for $n\ge m'$
\[
S_{n+1}\ge S_n\Bigl(1+\frac3n\Bigr)\ge S_n\Bigl(\frac{n+1}n\Bigr)^2,
\]
so $S_n\ge S_{m'}(n/m')^2$, which contradicts $S_n\le Dn$. Therefore $S_n=0$ for all $n$, and Lemma~\ref{lem:zeroS}
completes the proof.
\end{proof}

\begin{remark}\label{rem:dim}
Two features of the plane are used: $|B_n|$ grows linearly, and $\sum_n1/n$ diverges. Without \eqref{eq:odds}, the
truncation in \eqref{eq:tail} needs $\lambda$ of order $\log n$ to make $|B_n|\lambda e^{-\lambda}$ small, and the
resulting factor $\log n$ in \eqref{eq:planar} breaks the argument; this is the obstruction of
\cite[Remark~2.16]{MST25}. On $\Z^d$ with $d\ge3$ the same steps give $S_n^2\le Cn^{d-1}(S_{n+1}-S_n)$ and
$S_n\le Dn^{d-1}$, which are consistent with a bounded positive sequence $S$, so no conclusion follows. The comparison
with \cite{HS77}, where the same dichotomy between $d\le2$ and $d\ge3$ appears for rates bounded away from zero, is
that here the bound \eqref{eq:odds} on the measure, not a bound on the rates, keeps the boundary terms of
\eqref{eq:balance} of order $|B_n|$ times the boundary currents.
\end{remark}

\begin{proposition}\label{prop:stripbalance}
On a strip $\Z\times H$ with $H$ finite, every stationary measure satisfies detailed balance, for all constraints
of uniform finite range, possibly depending on the site, that do not depend on the spin they constrain.
\end{proposition}

\begin{proof}
Let $V_n=[-nr,nr]\times H$. Then every site of $V_{n-1}$ is interior to $V_n$ and $|B_n|=2r|H|$ for all $n$. By
Lemma~\ref{lem:shell}, $S_n$ is nondecreasing and bounded, so $S_{n+1}-S_n\to0$ and, by \eqref{eq:current},
$\sum_{x\in B_n}b_{V_n}(x)\to0$. By \eqref{eq:balance} and \eqref{eq:tail}, for every $\lambda\ge1$,
\[
S_n\le\lambda\sum_{x\in B_n}b_{V_n}(x)+4Mr|H|\lambda e^{-\lambda}.
\]
Letting $n\to\infty$ and then $\lambda\to\infty$ gives $\lim_nS_n=0$. As $S$ is nonnegative and nondecreasing,
$S_n=0$ for every $n$, and Lemma~\ref{lem:zeroS} applies.
\end{proof}

\section{Identification by internal spanning}\label{sec:ident}

From now on the constraints are monotone, given by \eqref{eq:update} on $\Z^2$ or of the form in
Proposition~\ref{prop:strip} on a strip.

\begin{lemma}\label{lem:connect}
Let $\mu$ satisfy detailed balance and let $V$ be finite. Then $\rho_V=\mu_V/\pi_V$ is constant on the configurations
of $\mathcal G_V$. Consequently, if $\mu(\mathcal G_V)>0$, then $\mu(\eta_V\in\cdot\mid\mathcal
G_V)=\pi(\eta_V\in\cdot\mid\mathcal G_V)$, and for every event $F$ depending only on $\eta_V$,
\begin{equation}\label{eq:compare}
|\mu(F)-\pi(F)|\le2\bigl(1-\mu(\mathcal G_V)\bigr)+2\bigl(1-\pi(\mathcal G_V)\bigr).
\end{equation}
\end{lemma}

\begin{proof}
Let $\zeta\in\mathcal G_V$. Emptying filled sites of $V$ in the order of bootstrap percolation inside $V$ with filled
exterior leads from $\zeta$ to the configuration $0_V$ in which every site of $V$ is empty. At each step the emptied
site $x$ has its constraint satisfied with every site outside $V$ filled, hence, by monotonicity, with any spins
outside $V$. Reversing such a path for another $\zeta'\in\mathcal G_V$ leads from $0_V$ to $\zeta'$; each reversed step
fills a site whose constraint is satisfied, because the constraint does not depend on the spin at that site. So any two
configurations of $\mathcal G_V$ are joined by a path of single flips at sites $x\in V$ whose constraint is satisfied
whatever the spins outside $V$.

For such a flip $\zeta\to\zeta^x$, let $W\supset V$ be finite and contain the sites on which $c_x$ depends. Then
$c_x(\xi)=1$ for every extension $\xi$ of $\zeta$ to $W$, and \eqref{eq:db} gives $r_x(\xi)\mu_W(\xi)=r_x(\xi^x)\mu_W(\xi^x)$.
Summing over the extensions gives $r_x(\zeta)\mu_V(\zeta)=r_x(\zeta^x)\mu_V(\zeta^x)$, that is
$\rho_V(\zeta)=\rho_V(\zeta^x)$, with both sides of the first identity zero or both positive. Hence $\rho_V$ is constant on
$\mathcal G_V$, which gives the identity of conditional laws. For \eqref{eq:compare}, for every probability $\nu$ with
$\nu(\mathcal G_V)>0$ we have $|\nu(F)-\nu(F\mid\mathcal G_V)|\le\nu(\mathcal G_V^c)+\nu(F\cap\mathcal
G_V)(1/\nu(\mathcal G_V)-1)\le2(1-\nu(\mathcal G_V))$; apply this to $\mu$ and to $\pi$. If $\mu(\mathcal G_V)=0$ the
bound \eqref{eq:compare} is trivial.
\end{proof}

\begin{lemma}\label{lem:domination}
If $\mu$ satisfies \eqref{eq:odds}, then $\mu(F)\ge\pi_\varepsilon(F)$ for every event $F$ that depends on finitely
many spins and is increasing in the set of empty sites.
\end{lemma}

\begin{proof}
Let $F$ depend on the spins at $x_1,\dots,x_N$. Let $U_1,\dots,U_N$ be independent uniform variables on $[0,1]$ and
define $\eta_{x_k}=0$ if $U_k\le\mu(\eta_{x_k}=0\mid\eta_{x_1},\dots,\eta_{x_{k-1}})$ and $\eta_{x_k}=1$ otherwise, the
conditional probability being evaluated at the values already constructed. The resulting spins have law
$\mu$. Each conditional probability is a conditional probability given $\eta_{V\setminus\{x_k\}}$ with
$V=\{x_1,\dots,x_k\}$ and is therefore at least $\varepsilon$ by \eqref{eq:odds}. Hence the set $\{x_k:U_k\le\varepsilon\}$, which has
law $\pi_\varepsilon$, is contained in the set of empty sites, and $F$ is increasing.
\end{proof}

\begin{proof}[Proof of Theorem~\ref{thm:identify}]
By Theorem~\ref{thm:balance}, $\mu$ satisfies detailed balance. The event $\mathcal G_{V_n}$ is increasing in the set of
empty sites, so Lemma~\ref{lem:domination} and the hypothesis give $\mu(\mathcal G_{V_n})\ge\pi_\varepsilon(\mathcal
G_{V_n})\to1$, and $\pi(\mathcal G_{V_n})=\pi_q(\mathcal G_{V_n})\to1$. By \eqref{eq:compare}, $\mu$ and $\pi$ agree on
every event depending on finitely many spins, so $\mu=\pi$.

For FA-1f, any empty site of the connected set $V_n$ empties all of $V_n$ by bootstrap percolation along paths, so
$\mathcal G_{V_n}=\{\eta_{V_n}\ne1_{V_n}\}$ and $\pi_a(\mathcal G_{V_n})=1-(1-a)^{|V_n|}\to1$. For FA-2f the hypothesis
is Lemma~\ref{lem:fa2f} below, a classical fact \cite{vE87,AL88}.
\end{proof}

\begin{proof}[Proof of Corollary~\ref{cor:fa1f}]
Both $\mu$ and $\delta_{\mathbf 1}$ are stationary, so $\nu$ is stationary; by Theorem~\ref{thm:identify}, $\nu=\pi$.
\end{proof}

\begin{proof}[Proof of Proposition~\ref{prop:strip}]
Let $\lambda=\mu(\{\mathbf 1\})$. If $\lambda=1$ there is nothing to prove. Otherwise
$\nu=(\mu-\lambda\delta_{\mathbf 1})/(1-\lambda)$ is a stationary probability measure with $\nu(\{\mathbf 1\})=0$, and
by Proposition~\ref{prop:stripbalance} it satisfies detailed balance. Since every FA-1f move is allowed and $V_n$ is
connected, a single empty site in $V_n$ empties $V_n$, so $\mathcal G_{V_n}=\{\eta_{V_n}\ne1_{V_n}\}$. These events
increase to $\{\mathbf 1\}^c$ along the increasing $V_n$, so $\nu(\mathcal G_{V_n})\to\nu(\{\mathbf 1\}^c)=1$, and
$\pi(\mathcal G_{V_n})=1-p^{|V_n|}\to1$. By Lemma~\ref{lem:connect}, $\nu=\pi$.
\end{proof}

\section{Internal spanning for FA-2f}\label{sec:fa2f}

The following lemma is classical; van Enter \cite{vE87} proved that two-neighbour bootstrap percolation on $\Z^2$
empties the whole plane at every positive density, and Aizenman and Lebowitz \cite{AL88} and Holroyd \cite{Hol03}
determined the scale at which boxes become internally spanned. We include a short proof of the form used above.

\begin{lemma}\label{lem:fa2f}
For FA-2f and every $a\in(0,1)$, $\pi_a(\mathcal G_{V_n})\to1$ as $n\to\infty$.
\end{lemma}

\begin{proof}
By translation it suffices to treat $Q_N=[1,N]^2$. Put $t=1-a$, fix $C>3/|\log t|$, and let $L=\lceil C\log N\rceil$
and $m=\lceil(\log N)^{1/3}\rceil$, so that $m<L<N$ for large $N$. Two facts about two-neighbour bootstrap percolation
are used. If a $k\times j$ rectangle $R\subset Q_N$ is empty and a segment $T$ of $k$ (or $j$) sites adjacent to a
side of $R$ contains an empty site $y$, then all of $T$ empties: the neighbours of $y$ in $T$ have two empty neighbours,
one in $R$ and $y$, and the emptying propagates along $T$. Every site involved lies in $Q_N$, so the exterior may stay
filled.

Partition a subsquare of $Q_N$ into $J=\lfloor N/L\rfloor^2$ disjoint $L\times L$ blocks. Call a block good if its
bottom-left $m\times m$ square is initially empty and, for $k=m,\dots,L-1$, the column segment of length $k$ to the
right of the bottom-left $k\times k$ square and the row segment of length $k+1$ above the resulting $k\times(k+1)$
rectangle each contain an initial empty site. These segments and the seed square are disjoint, so by the first fact a
good block is internally emptied, and
\[
\pi_a(\text{good})\ge a^{m^2}\prod_{k=m}^{L-1}(1-t^k)(1-t^{k+1})\ge a^{m^2}\Bigl(1-\frac{2t^m}{1-t}\Bigr)\ge\frac{a^{m^2}}2
\]
for large $N$. Blocks are independent, so the probability that no block is good is at most
$\exp(-Ja^{m^2}/2)$, which tends to $0$ because $\log J=2\log N-O(\log\log N)$ and $m^2=O((\log N)^{2/3})$.

Let $\mathcal S$ be the event that some horizontal or vertical segment of $L$ sites in $Q_N$ is initially filled. There
are at most $2N^2$ such segments, so $\pi_a(\mathcal S)\le2N^2t^L\le2N^{2-C|\log t|}\to0$. Suppose some block is good
and $\mathcal S$ fails. The emptied block has height $L$; each column segment of length $L$ next to it contains an
empty site, so by the first fact the block grows column by column into an empty band of height $L$ across $Q_N$. Each
row of $Q_N$ next to the band contains an empty site, since it contains a segment of length $L$, so the band grows row
by row until $Q_N$ is empty. Hence $\pi_a(\mathcal G_{Q_N})\ge1-\exp(-Ja^{m^2}/2)-2N^2t^L\to1$.
\end{proof}

\begin{remark}\label{rem:finite}
In a finite box with fixed exterior, legal flips preserve the bootstrap limit, and any two configurations with the
same limit are joined by legal flips, as in the proof of Lemma~\ref{lem:connect}. So the communicating classes of the
finite chain are the bootstrap fibers, the restriction of $\pi$ to each fiber is reversible, and the stationary
measures of the finite chain are exactly the mixtures of $\pi$ conditioned on the fibers. The difficulty in infinite
volume is the passage from stationarity to detailed balance, which is what Theorem~\ref{thm:balance} supplies under
\eqref{eq:odds}.
\end{remark}

\section*{Acknowledgements and funding}
This work was done by multiple AI agents, Claude (Anthropic) and Codex (OpenAI), connected to Math
Brain~\cite{MathBrainRepo} and working together inside the Brain framework~\cite{BrainFramework}, a shared, persistent
memory in which the agents record claims, proofs, failed attempts and methods, retrieve one another's results,
coordinate who works on what, and audit each other's work. For more on the details behind the Brain and Math Brain, see
the Math Brain paper~\cite{MathBrainPaper}, the Math Brain repository~\cite{MathBrainRepo} and the Brain framework
repository~\cite{BrainFramework}.

The research presented in this paper was supported by the European Research Council (ERC) under the European Union's
Horizon Europe research and innovation programme (grant agreement No.~101041711), by the Simons Foundation as part of
the Collaboration on the Mathematical and Scientific Foundations of Deep Learning, by Heights Labs, by Convex Nexus
Capital, by the Israel Science Foundation (grant number 2258/19), and by the Israel Science Foundation (ISF Grant
4101/25). Computational resources were provided in part by Google Cloud through its education credits programme (credit
530550103).

\section*{Declaration of competing interest}
The author holds equity in Heights Labs and in Convex Nexus Capital and has a management role at Convex Nexus Capital.
These entities are included in the funding declaration above.

\section*{Declaration of generative AI and AI-assisted technologies}
A Codex agent (OpenAI) proved Theorems~\ref{thm:balance} and~\ref{thm:identify}, Proposition~\ref{prop:strip} and
Lemma~\ref{lem:fa2f}; a second Codex agent checked the proofs as a referee; Claude agents (Anthropic) checked the steps
and prepared the manuscript. The author directed the work and is responsible for the content of the paper and for
decisions concerning its dissemination.

\section*{Code availability}
A script that checks Remark~\ref{rem:finite} exactly for FA-1f and FA-2f on boxes of up to eight sites, with filled,
empty and periodic boundary conditions and $q=1/3$ and $5/7$, is in the repository
\begin{center}\url{https://github.com/yspennstate/kcm-planar-stationary-bounded-odds}.\end{center}
For each of the $40$ cases it verifies in rational arithmetic that $\pi$ restricted to each bootstrap fiber is stationary
and reversible, and that the dimension of the space of stationary vectors equals the number of fibers.

\appendix
\section{Notation}\label{app:notation}

\begin{center}
\begin{tabular}{p{3.3cm}p{7.6cm}p{3.2cm}}
\toprule
Symbol & Meaning & First appears\\
\midrule
$G$, $\Omega$, $\eta$, $\eta_x$ & lattice or strip, configurations, a configuration, the spin at $x$ ($0$ empty, $1$ filled) & Section~\ref{sec:intro}\\
$q$, $p$, $\pi$, $\pi_a$ & density of empty sites, $p=1-q$, product measures & Section~\ref{sec:intro}\\
$c_x$, $r$ & constraint at $x$, its range & Section~\ref{sec:intro}\\
$r_x$, $\eta^x$, $L$, $P_t$ & flip rate, flipped configuration, generator, semigroup & \eqref{eq:generator}\\
$\cU$ & update family & \eqref{eq:update}\\
$\BP$, $E$, $\mathcal G_V$ & bootstrap limit, stable configurations, internal spanning of $V$ & Section~\ref{sec:intro}\\
$\mathbf 1$, $1_V$, $0_V$ & all sites filled, all sites of $V$ filled, all sites of $V$ empty & Section~\ref{sec:intro}, Lemma~\ref{lem:connect}\\
$\mathbb I$ & indicator & \eqref{eq:update}\\
$H$ & finite connected graph, the cross-section of a strip & Section~\ref{sec:intro}\\
$f$, $F$ & local function, event depending on finitely many spins & Section~\ref{sec:intro}, Lemma~\ref{lem:connect}\\
$U$, $u$ & element of an update family, a site of it & \eqref{eq:update}\\
$\beta$, $\mu^*$ & stable configuration, mixing measure & Section~\ref{sec:intro}\\
$\zeta$, $\xi$ & configurations of a finite set & \eqref{eq:db}\\
$\phi$ & $(u,v)\mapsto(u-v)\log(u/v)$ & Lemma~\ref{lem:abc}\\
$U_k$ & independent uniform variables & Lemma~\ref{lem:domination}\\
$R$, $T$, $\mathcal S$ & empty rectangle, adjacent segment, a filled segment of length $L$ exists & Lemma~\ref{lem:fa2f}\\
$\mu_W$, $\mu$, $\nu$ & marginal on $W$, stationary measures & \eqref{eq:db}\\
$\varepsilon$ & constant in the bounded odds condition & \eqref{eq:odds}\\
$a_V$, $b_V$, $A_V$ & integrated flip rates, their imbalance, entropy production & Section~\ref{sec:entropy}\\
$\rho_V$, $\ell_V$ & density of $\mu_V$ with respect to $\pi_V$, log ratio across a flip & Section~\ref{sec:entropy}\\
$M$, $\lambda$ & $\max(p/q,q/p)$, truncation level & Lemma~\ref{lem:balance}\\
$V_n$, $I_n$, $B_n$, $S_n$ & boxes, interiors, shells, entropy production in $V_n$ & Section~\ref{sec:entropy}\\
$K$, $C$, $D$ & constants of the planar estimate & Section~\ref{sec:proofs}\\
$Q_N$, $L$, $m$, $J$, $t$ & square, block size, seed size, number of blocks, $1-a$ & Lemma~\ref{lem:fa2f}\\
\bottomrule
\end{tabular}
\end{center}

\begin{thebibliography}{99}
\bibitem{AL88} M.~Aizenman and J.~L.~Lebowitz, Metastability effects in bootstrap percolation, \emph{J.\ Phys.\ A}
\textbf{21} (1988), no.~19, 3801--3813. \doi{10.1088/0305-4470/21/19/017}
\bibitem{BCMRT13} O.~Blondel, N.~Cancrini, F.~Martinelli, C.~Roberto and C.~Toninelli, Fredrickson--Andersen one spin
facilitated model out of equilibrium, \emph{Markov Process.\ Related Fields} \textbf{19} (2013), no.~3, 383--406.
\arxiv{1205.4584}
\bibitem{CMRT08} N.~Cancrini, F.~Martinelli, C.~Roberto and C.~Toninelli, Kinetically constrained spin models,
\emph{Probab.\ Theory Related Fields} \textbf{140} (2008), no.~3--4, 459--504. \doi{10.1007/s00440-007-0072-3}
\bibitem{FA84} G.~H.~Fredrickson and H.~C.~Andersen, Kinetic Ising model of the glass transition, \emph{Phys.\ Rev.\
Lett.} \textbf{53} (1984), no.~13, 1244--1247. \doi{10.1103/PhysRevLett.53.1244}
\bibitem{HT25} I.~Hartarsky and C.~Toninelli, \emph{Kinetically Constrained Models}, SpringerBriefs in Mathematical
Physics, Springer, Cham, 2025. \doi{10.1007/978-3-031-93115-4}
\bibitem{Hol71} R.~Holley, Free energy in a Markovian model of a lattice spin system, \emph{Comm.\ Math.\ Phys.}
\textbf{23} (1971), no.~2, 87--99. \doi{10.1007/BF01877751}
\bibitem{HS77} R.~Holley and D.~W.~Stroock, In one and two dimensions, every stationary measure for a stochastic Ising
model is a Gibbs state, \emph{Comm.\ Math.\ Phys.} \textbf{55} (1977), no.~1, 37--45. \doi{10.1007/BF01613147}
\bibitem{Hol03} A.~E.~Holroyd, Sharp metastability threshold for two-dimensional bootstrap percolation, \emph{Probab.\
Theory Related Fields} \textbf{125} (2003), no.~2, 195--224. \doi{10.1007/s00440-002-0239-x}
\bibitem{JK25} B.~Jahnel and J.~K\"oppl, On the long-time behaviour of reversible interacting particle systems in one
and two dimensions, \emph{Probab.\ Math.\ Phys.} \textbf{6} (2025), no.~2, 479--503. \doi{10.2140/pmp.2025.6.479}
\bibitem{Mar26} F.~Martinelli, Kinetically constrained spin models: a short survey, in \emph{Proceedings of the
International Congress of Mathematicians 2026}, vol.~5, 497--515. \doi{10.1137/25M180425X}
\bibitem{MST25} F.~Martinelli, A.~Shapira and C.~Toninelli, Long time behaviour of one facilitated kinetically
constrained models: results and open problems, preprint, 2025. \arxiv{2510.20461}
\bibitem{MOP77} J.~Moulin Ollagnier and D.~Pinchon, Free energy in spin-flip processes is non-increasing, \emph{Comm.\
Math.\ Phys.} \textbf{55} (1977), no.~1, 29--35. \doi{10.1007/BF01613146}
\bibitem{MV19} T.~Mountford and G.~Valle, Exponential convergence for the Fredrickson--Andersen one-spin facilitated
model, \emph{J.\ Theoret.\ Probab.} \textbf{32} (2019), no.~1, 282--302. \doi{10.1007/s10959-018-0812-6}
\bibitem{vE87} A.~C.~D.~van Enter, Proof of Straley's argument for bootstrap percolation, \emph{J.\ Stat.\ Phys.}
\textbf{48} (1987), no.~3--4, 943--945. \doi{10.1007/BF01019705}
\bibitem{MathBrainRepo} Y.~Shmalo, \emph{Math Brain: public snapshot, code and tests}, 2026,
\url{https://github.com/yspennstate/math-brain-public}.
\bibitem{MathBrainPaper} Y.~Shmalo, \emph{Math Brain: shared, auditable memory for multi-agent mathematical research},
October 2026, \url{https://github.com/yspennstate/math-brain-public/tree/main/papers/math-brain-2026-10-24}.
\bibitem{BrainFramework} Y.~Shmalo, \emph{Brain Framework}, public software repository, 2026,
\url{https://github.com/yspennstate/brain-framework-public}.
\end{thebibliography}
\end{document}
